\subsection{Models}

We train five model families targeting $r^e_{i,t}$, sharing the same temporal split and feature sets.

\textbf{XGBoost variants (four models).}
\begin{enumerate}[noitemsep]
    \item \textbf{JKP only}: 153 lagged factor features.
    \item \textbf{FinBERT only}: 12 sentiment features (Table~\ref{tab:features}). \textit{Primary model.}
    \item \textbf{JKP + FinBERT}: all 165 features.
    \item \textbf{JKP + FinBERT $\times$10}: same 165 features, FinBERT columns weighted $10\times$ via \texttt{feature\_weights} in \texttt{DMatrix}.
\end{enumerate}
Shared hyperparameters: 500 estimators, early stopping at 30 rounds on validation RMSE, learning rate 0.05, max depth 4, subsample 0.8, $L_2$ regularisation 1.0, \texttt{random\_state=42}.

\textbf{Deep learning models (PyTorch).} The project requires at least one deep learning approach; the four models below additionally serve as cross-model robustness checks:
\begin{enumerate}[noitemsep]
    \item \textbf{MLP\textsubscript{FinBERT}}: fully-connected MLP on the 12 FinBERT features. Architecture and training details in Appendix~B.
    \item \textbf{MLP\textsubscript{JKP+FinBERT}}: MLP on all 165 features (neural-network equivalent of combined XGBoost).
    \item \textbf{LSTM\textsubscript{FinBERT}}: two-layer LSTM on the temporal sequence of the last 8 quarterly calls (Appendix~B).
    \item \textbf{Attention\textsubscript{FinBERT}}: Pre-LN Transformer encoder on the same 8-step sequence (Appendix~B).
\end{enumerate}
All DL models use Adam ($\ell_r = 10^{-3}$), early stopping on validation Rank IC (patience 15, max 150 epochs), and Apple M4 MPS backend. MLP\textsubscript{FinBERT} operates on the same 12 features as the primary XGBoost model, enabling a direct comparison; XGBoost outperforming it is an empirical result, not a weakness of the DL approach.

\textbf{Simple linear benchmark.} We additionally run a Ridge regression ($\alpha=1.0$) on the full 12 FinBERT features to test whether the signal requires non-linear modelling.

\subsection{Evaluation Metrics}

\begin{itemize}[noitemsep]
    \item \textbf{Rank IC}: mean monthly Spearman correlation between predicted and realised excess returns, computed cross-sectionally within each month. This is the primary evaluation metric, as the strategy profits from within-month stock rankings rather than market-direction calls \cite{grinold1999}. Global Rank IC (Spearman across all stock-month pairs) mixes time-series and cross-sectional variation and is reported for context only.
    \item \textbf{IC $t$-statistic}: $\bar{\text{IC}} / (\sigma_\text{IC} / \sqrt{T})$ over $T$ monthly ICs.
\end{itemize}

\subsection{Investment Strategies and Transaction Costs}

Equal-weighted long-short decile strategy: long top 10\%, short bottom 10\%, rebalanced monthly. Transaction costs of \textbf{0.2\% one-way} (0.4\% round-trip) are applied on monthly turnover. We report annualised gross and net returns, Sharpe ratio, OLS alpha versus S\&P500, and average monthly turnover.
